\section{The Three Readings Compared}
\label{sec:comparison}

The three readings hear the same ten texts and agree on much. All three read the
formulary under one sovereign God who is refuge: the Introit's will that no one
resists and the Gradual's God who is before the mountains stand behind each of
them. All three hear the Gospel's mercy and the Communion's plea for judgment
together, and none lets the plea undo the command to forgive. All three draw on
St John Chrysostom, who expounds both the Epistle and the Gospel, and on St Thomas
Aquinas, whose distinctions serve each in turn: the will and the power to harm in
the first, the sins said to return through ingratitude in the second, and the
reckoning that ends at the judgment in the second and third. The second and third
rest on St Augustine, and the first and third draw on St Gregory the Great. Where
they differ is in the question each asks of the texts,
the element each hears first, and the answer each gives.

\begin{comparisontable}{}{The war God permits and bounds}{The king's reckoning and the forgiven debtor}{The Church's exodus toward the judgment}
Question & Who is the Christian's enemy, and how can his work stand within the
will no one resists? & What does the king's mercy require of the forgiven, and
can a pardon return? & When will the judgment the Communion asks for come, and
where is the people meanwhile?\\
Heard first & Introit and Offertory, with the Epistle & Gospel & Communion, with
the Alleluia, the Epistle's girded loins and the Postcommunion\\
Decisive inference & The devil asks and God permits, bridling what he looses
(Gregory); the will to harm is the devil's, the power God's (Thomas); the spirits
that move persecutors are the real enemy (Chrysostom) & The receiver revokes the
gift by his cruelty (Chrysostom); God calls back what he remitted (Augustine);
forgiven sins return through ingratitude, not as such (Thomas) & The plea is the
martyrs', answered with long-suffering (Augustine); the question asks for the end,
not for the length of a life (Cassiodorus, Bellarmine); the Christian life is an
Exodus, eaten girded (Chrysostom)\\
The Communion's persecutors & The devil and his ministers, against whom the singer
asks the power of Lk 10:19 (Hilary, on the sung half; his ``now'' drawn from the
unsung question) & The dominion of sin, whose overthrow avenges the martyrs, and
men who are to be forgiven (Augustine) & The Church's persecutors, judged at
the end (Augustine); the wicked who persecute and tempt (Bellarmine); the devil
and his ministers (Cassiodorus)\\
What it asks of the hearer & Stand, and turn the wrath from the brother to the
devil & Forgive every debt from the heart, and correct without rancour & Live as
on a march, and forgive while God's patience waits\\
\end{comparisontable}

\subsection{Where the readings meet}

The first and second readings meet at a sentence of Chrysostom's. In his homily on
the Epistle he tells the man who has been wronged to pour all his wrath on the
devil; in his homily on the parable he asks why a man who has seen the devil
become a cause of crowns should fear another man as an enemy. The first reading
takes the first half of that thought: the enemy is the adversary, not the
brother. The second takes the other half: the brother is the fellow servant whose
small debt the forgiven man must forgive. The two are complementary, and they are
distinct. The first reading's argument stands on the Introit, the Epistle, the
Offertory and the Communion, and would stand without the parable; the second's
stands on the parable, and would stand without the Offertory's permission. The
first answers a question about God's government; the second, a question about
the forgiven man's heart.

The second and third readings meet at Augustine's sermon on anger, where the
martyrs not yet avenged and the command ``from the heart forgive'' stand in one
sentence. The second reading uses it to show that the Communion's plea does not
undo the Gospel; the third, to show that the waiting it imposes is the time of
God's patience with persecutors. Both hold that the Church that asks for judgment
forgives in the meantime. They differ in where they place the weight: the second
on the debt and its condition, the third on the time and its end.

\subsection{Where they differ}

The sharpest difference lies in the Communion, and it is a difference between
Fathers, not only between readings. St Hilary asks for a power given in this life
to crush the devil and his ministers, and says that this is not the judgment of
the resurrection, which he affirms. St Augustine hears the martyrs' cry for the
final judgment, when vengeance falls on the Church's persecutors; St Robert
Bellarmine, a plea for the condemnation of the wicked who persecute and tempt;
Cassiodorus joins Hilary's persecutors, the devil with his ministers, to the
judgment at the end of the world. They differ in the thing asked and in the
persecutors named, and also in the time. Each draws
his time from the question the Communion does not sing, ``How many are the days of
thy servant?'': Hilary hears in it a man who knows his days are few; Augustine,
the body of Christ asking about its days in the world; Cassiodorus and Bellarmine,
no question about the length of a life at all. The exposition printed under St
Athanasius's name holds both answers: on that unsung question it asks that Satan
be crushed at once, and on the sung cry it longs to see the time of the divine
tribunal. The answers are compatible, though they are not the same. Hilary
affirms the last judgment and the devil's judgment at the resurrection; only he
and Cassiodorus name the persecutors as the devil with his ministers, while
Augustine's and Bellarmine's persecutors are men, whom the first reading, with St
Thomas, calls instruments moved by the devil. Hilary's identification of the
persecutors stands on the sung words, his ``now'' on the unsung question;
Augustine's martyrs' plea answered by long-suffering stands on the sung words, his
``until the judgment'' on the order of the two questions. The first reading
follows Hilary's identification and the power he asks for; the third follows the
end-directed answer. The difference is real, and it divides no doctrine.

A second difference lies in the Gospel's sentence. The second reading keeps three
formulations side by side: Chrysostom's servant who revokes his own pardon,
Augustine's God who calls back what he remitted, and Thomas's sins that do not
return as such but are said to return through ingratitude. The first two are
complementary in their own grammar, the one looking at the cause and the other at
the judge; Thomas's is narrower than either, and he names Origen as holding the
stronger view he denies. The third reading reads the same king's reckoning toward
its end at the judgment, with Thomas, and does not need to decide the question of
return.

A third lies in the armour. St Jerome makes the armour of God Christ himself;
Thomas makes it the gifts and virtues; the continuation of \work{The Liturgical
Year} makes it first Christ's and then ours. The first reading holds Jerome's
identification and Thomas's together as the Christological and the moral sides of
one gift. The third reading draws from the same Epistle not the armour but the
girded loins, which Chrysostom reads as the girding of the Passover.

\subsection{What each reading meets in the texts}

Each reading meets a demanding feature of the formulary that the others leave in
the background. The first takes up the Introit's absolute sentence beside the
Offertory's \latin{petiit} and \latin{data est \dots\ a Domino}, the hardest
theological knot in the Mass, and answers it with Gregory's two wills and Thomas's
distinction of will and power. The second takes up the Gospel's reversed pardon
and its \latin{universum debitum}, and answers it without softening the sentence
or letting it unsettle the efficacy of God's forgiveness. All three take up the
plea for judgment sung after a Gospel of forgiveness. The first answers it from
Hilary, for whom the persecutors are the devil and his ministers; the second from
Augustine, for whom the martyrs plead against the dominion of sin and not against
the men; the third, again from Augustine, by the time the plea opens, the delay of
the judgment in which God's patience waits for the persecutors' conversion. Read together, the
three readings show one Mass in which the God who bounds the adversary is the king
who forgives the debt, and the people that stands in the evil day is the people
that goes out of Egypt toward his judgment.
